% ---------------------------------------------------------------------
%  §3. Datos y metodología
% ---------------------------------------------------------------------
\section{Datos y metodología}

\textbf{Estructura.} La base contiene \textbf{1.743 registros} y \textbf{11 variables} sobre estudiantes de
18 a 24 años de \textbf{110 países}, con una distribución de género pareja (883 varones y 860 mujeres). Las
variables centrales son el \textbf{uso diario de redes} (de 1,5 a 8,5 horas; media 4,98), las \textbf{horas
de sueño} por noche (de 3,8 a 9,6; media 6,78), el \textbf{puntaje de salud mental} (escala de 4 a 9, más
alto = mejor; media 6,21), el \textbf{impacto académico percibido} (si el estudiante \emph{declara} que las
redes afectan su rendimiento; es una percepción y no una medición de sus notas), la \textbf{plataforma más
usada} (12 redes, con Instagram, 550, y TikTok, 350, concentrando más de la mitad) y el \textbf{estado
sentimental} (soltero/a, en pareja o complicado), además de país, género, edad y nivel académico.

\textbf{Calidad de los datos.} No se detectaron valores atípicos, pero sí dos problemas relevantes. El
estado sentimental es la única variable incompleta, con 1.038 valores faltantes (59,55\,\%). Además,
\textbf{705 estudiantes} aparecían cargados dos veces —una sin estado sentimental y otra con ese dato—,
por lo que a simple vista parecían filas distintas. El identificador de estudiante tampoco es único: hay
números compartidos por personas diferentes (distinto género, país y edad), que se conservaron.

\textbf{Transformaciones.}
\begin{itemize}
  \item \textbf{Nivel de desarrollo del país:} según la clasificación por ingreso del Banco Mundial
  \cite{bancomundial}, las economías de ingreso alto se agruparon como \emph{Desarrollado} y el resto como
  \emph{En desarrollo} (1.103 y 638 estudiantes). Ciudad del Vaticano (2 registros) no figura en esa
  clasificación y quedó excluida de esta comparación.
  \item \textbf{Eliminación de duplicados:} para el análisis por plataforma y género se dejó una sola copia
  de cada estudiante repetido, con lo que la base quedó en \textbf{1.038 estudiantes}.
  \item \textbf{Plataformas minoritarias:} WeChat, LINE, KakaoTalk y VKontakte, con muy pocos casos, se
  agruparon en la categoría \emph{Otras} para evitar promedios inestables.
  \item \textbf{Niveles de uso:} \emph{Bajo} (hasta 4 h, 278 estudiantes), \emph{Medio} (entre 4 y 6 h,
  493) y \emph{Alto} (más de 6 h, 267), para comparar plataformas entre estudiantes con un uso similar.
  \item \textbf{Impacto académico como proporción:} la respuesta Sí/No se convirtió en un indicador cuyo
  promedio equivale al porcentaje de estudiantes que se sienten afectados.
\end{itemize}

Las preguntas 1 y 2 se analizan sobre los 1.743 registros (al repetir los cálculos principales sin
duplicados los resultados cambian solo en los decimales y las conclusiones se mantienen), la pregunta 3
sobre los 1.038 estudiantes sin duplicados y la pregunta 4 sobre los 705 registros con estado sentimental
completo.

\textbf{Enfoque y herramientas.} El análisis es \textbf{descriptivo}: se comparan medias, medianas y
proporciones entre grupos, se calculan correlaciones y rectas de tendencia, y se construyen histogramas,
diagramas de caja, gráficos de barras y de dispersión. El procesamiento se realizó en Python (pandas,
NumPy, Matplotlib y seaborn) en Google Colab \cite{mckinney}. Los resultados describen asociaciones dentro
de la muestra y no permiten establecer relaciones causales.
